\documentclass[11pt,a4paper]{article}

\usepackage[T1]{fontenc}
\usepackage[utf8]{inputenc}
\usepackage[french]{babel}

\usepackage[a4paper,margin=2.2cm]{geometry}
\usepackage{amsmath,amssymb,amsfonts,amsthm}
\usepackage{mathtools}
\usepackage{bm}
\usepackage{graphicx}
\usepackage{xcolor}
\usepackage{hyperref}
\usepackage{enumitem}
\usepackage{booktabs}
\usepackage{array}
\usepackage{tikz}
\usepackage{fancyhdr}

\hypersetup{
    colorlinks=true,
    linkcolor=blue!60!black,
    urlcolor=blue!60!black,
    citecolor=blue!60!black,
    pdftitle={MedVision -- Fondements mathématiques et algorithmiques},
    pdfauthor={Yann Smatti}
}

\setlength{\parindent}{0pt}
\setlength{\parskip}{0.6em}

\pagestyle{fancy}
\fancyhf{}
\lhead{MedVision — Fondements mathématiques}
\rhead{\thepage}

\title{\textbf{MedVision}\\
Fondements Mathématiques et Algorithmiques du Pipeline\\
\large Analyse détaillée avec interprétations intuitives\\
pour étudiants en L2--L3 / CPGE}
\author{Yann Smatti}
\date{Mars 2026}

\begin{document}

\maketitle

\begin{center}
\textbf{Pipeline ML · TensorFlow · Computer Vision · Imagerie médicale}
\end{center}

\tableofcontents
\newpage

\section{Introduction et prérequis}

Ce document explique les fondements mathématiques du projet \textbf{MedVision},
un pipeline de vision par ordinateur appliqué à l'analyse d'images médicales.

L'idée générale du projet est la suivante :
\begin{enumerate}[label=\arabic*)]
    \item on charge des images médicales ;
    \item on les prétraite ;
    \item on les passe dans un réseau de neurones convolutif ;
    \item on entraîne le modèle à reconnaître une classe pathologique ou non ;
    \item on évalue rigoureusement ses performances.
\end{enumerate}

\textbf{Prérequis utiles :}
\begin{itemize}
    \item algèbre linéaire : vecteurs, matrices, produit scalaire ;
    \item analyse : dérivées partielles, gradient, règle de chaîne ;
    \item probabilités : variables aléatoires, espérance, variance ;
    \item un peu d’informatique scientifique.
\end{itemize}

\textbf{Philosophie du document :} chaque notion est présentée en trois niveaux :
\begin{itemize}
    \item l’idée intuitive ;
    \item la formulation mathématique ;
    \item le lien direct avec le projet MedVision.
\end{itemize}

\section{Représentation mathématique d'une image}

\subsection{Une image comme fonction discrète}

Une image numérique en niveaux de gris peut être vue comme une fonction

\[
I : \{1,\dots,H\} \times \{1,\dots,W\} \to [0,255]
\]

où :
\begin{itemize}
    \item $H$ désigne la hauteur ;
    \item $W$ désigne la largeur ;
    \item $I(i,j)$ est l’intensité lumineuse du pixel situé en position $(i,j)$.
\end{itemize}

Pour une image couleur RGB, on a en réalité trois canaux, donc

\[
I \in \mathbb{R}^{H \times W \times 3}.
\]

\textbf{Intuition.}  
Une image couleur est donc une sorte de \emph{tenseur d’ordre 3} :
hauteur, largeur, canal.

\subsection{Image comme vecteur de grande dimension}

Si l’on « déroule » l’image, on peut aussi la considérer comme un vecteur

\[
x \in \mathbb{R}^{3HW}.
\]

Par exemple, pour une image $224 \times 224 \times 3$ :

\[
3HW = 3 \times 224 \times 224 = 150\,528.
\]

Donc une seule image vit dans un espace de dimension 150\,528.

\textbf{Conséquence.}  
Le rôle du modèle est d’apprendre une fonction

\[
f : \mathbb{R}^{150528} \to \mathbb{R}^K,
\]

où $K$ est le nombre de classes.

\section{Prétraitement des images}

\subsection{Normalisation}

Les pixels d’une image sont souvent codés entre $0$ et $255$.
Pour stabiliser l’entraînement, on normalise :

\[
x' = \frac{x}{255}.
\]

Ainsi :

\[
x' \in [0,1].
\]

\textbf{Pourquoi ?}  
Parce que les réseaux de neurones apprennent mieux quand les données sont sur des échelles comparables.

\subsection{Redimensionnement}

Dans MedVision, les images sont redimensionnées en général à

\[
224 \times 224
\]

pour être compatibles avec des architectures pré-entraînées standards
comme ResNet ou EfficientNet.

\textbf{Interprétation.}  
On impose un format commun à toutes les images, exactement comme on met tous les vecteurs dans le même espace avant de faire de l’algèbre linéaire.

\subsection{Augmentation des données}

On applique souvent des transformations comme :
\begin{itemize}
    \item rotation ;
    \item symétrie horizontale ;
    \item zoom léger ;
    \item variation de contraste.
\end{itemize}

Mathématiquement, on applique une transformation

\[
T : \mathcal{X} \to \mathcal{X}
\]

à une image $x$, ce qui donne une nouvelle image

\[
\widetilde{x} = T(x).
\]

\textbf{Idée.}  
On ne change pas la nature médicale de l’image, mais on enrichit artificiellement le dataset.

\section{Convolution}

\subsection{Définition discrète}

La convolution discrète entre une image $I$ et un noyau $K$ est définie par

\[
(I * K)(i,j)=\sum_m \sum_n I(i-m,j-n)K(m,n).
\]

Dans la pratique, le noyau est petit, par exemple $3\times 3$ ou $5\times 5$.

\subsection{Interprétation intuitive}

Pour calculer la nouvelle valeur d’un pixel, on combine les valeurs des pixels voisins avec des poids donnés par le noyau.

Exemple de noyau détectant des contours verticaux :

\[
K=
\begin{pmatrix}
1 & 0 & -1 \\
1 & 0 & -1 \\
1 & 0 & -1
\end{pmatrix}.
\]

Si l’image contient un bord vertical, ce filtre réagit fortement.

\subsection{Lien avec l’analyse}

En analyse continue, la convolution est

\[
(f*g)(x)=\int_{\mathbb{R}} f(t)g(x-t)\,dt.
\]

En image, on passe simplement à une version discrète en deux dimensions.

\textbf{Idée importante.}  
La convolution agit comme un \emph{filtrage local} : elle extrait des motifs élémentaires comme
les bords, textures, petits contrastes.

\section{Le neurone artificiel}

Un neurone artificiel calcule

\[
y = \varphi(w^\top x + b),
\]

où :
\begin{itemize}
    \item $x$ est le vecteur d’entrée ;
    \item $w$ est le vecteur de poids ;
    \item $b$ est un biais ;
    \item $\varphi$ est une fonction d’activation.
\end{itemize}

\subsection{Interprétation géométrique}

L’équation

\[
w^\top x + b = 0
\]

définit un hyperplan.

Donc un neurone simple sépare l’espace en deux régions.

\textbf{Vision géométrique.}  
Un neurone décide de quel côté d’un hyperplan se trouve le point $x$.

\section{Réseaux de neurones convolutifs}

\subsection{Principe}

Un CNN empile plusieurs blocs :
\begin{itemize}
    \item convolution ;
    \item activation ;
    \item pooling ;
    \item couches denses finales.
\end{itemize}

Une couche convolutive calcule

\[
Y_{i,j,f}
=
\sum_{c=1}^{C}
\sum_{m=0}^{k-1}
\sum_{n=0}^{k-1}
W_{m,n,c,f}X_{i+m,j+n,c} + b_f.
\]

Ici :
\begin{itemize}
    \item $X$ est l’entrée ;
    \item $W$ est le filtre ;
    \item $f$ indexe la carte de caractéristiques obtenue.
\end{itemize}

\subsection{Pourquoi c’est plus efficace qu’une couche dense ?}

Une couche dense relie tous les pixels à tous les neurones.
Une convolution :
\begin{itemize}
    \item utilise des voisinages locaux ;
    \item partage les mêmes poids sur toute l’image.
\end{itemize}

Donc elle a beaucoup moins de paramètres.

\section{Fonctions d’activation}

\subsection{ReLU}

La fonction la plus classique est

\[
\mathrm{ReLU}(x)=\max(0,x).
\]

\textbf{Pourquoi ?}  
Elle introduit une non-linéarité indispensable. Sans non-linéarité, un réseau profond resterait essentiellement une transformation linéaire globale.

\subsection{Sigmoïde}

Pour une classification binaire, on utilise souvent :

\[
\sigma(x)=\frac{1}{1+e^{-x}}.
\]

Cette fonction renvoie une probabilité dans $[0,1]$.

\section{Pooling}

Le pooling sert à réduire la taille spatiale des représentations.

Exemple en max-pooling sur une fenêtre $2\times 2$ :

\[
\begin{pmatrix}
1 & 3 \\
2 & 0
\end{pmatrix}
\mapsto 3.
\]

\textbf{Interprétation.}  
On garde la réponse la plus forte localement. Cela rend le modèle plus robuste à de petits déplacements des motifs.

\section{Descente de gradient}

\subsection{Principe général}

On cherche à minimiser une fonction de coût $L(\theta)$.
La mise à jour de base est

\[
\theta_{t+1} = \theta_t - \eta \nabla L(\theta_t),
\]

où :
\begin{itemize}
    \item $\theta$ désigne les paramètres ;
    \item $\eta$ est le taux d’apprentissage ;
    \item $\nabla L$ est le gradient.
\end{itemize}

\subsection{Interprétation}

Le gradient pointe dans la direction de plus forte augmentation.
Donc, pour minimiser, on se déplace dans la direction opposée.

\textbf{Analogie.}  
C’est comme descendre une montagne en regardant la pente locale.

\section{Rétropropagation}

La rétropropagation repose sur la règle de chaîne.

Si

\[
f(x)=f_3(f_2(f_1(x))),
\]

alors

\[
\frac{df}{dx}
=
\frac{df_3}{df_2}
\cdot
\frac{df_2}{df_1}
\cdot
\frac{df_1}{dx}.
\]

\textbf{Sens.}  
On calcule comment une petite variation d’un paramètre en amont influence la perte finale.

C’est cette idée qui permet d’entraîner efficacement des réseaux profonds.

\section{Fonction de perte}

\subsection{Classification binaire : entropie croisée}

Pour une vérité terrain $y \in \{0,1\}$ et une probabilité prédite $p$ :

\[
L(y,p) = -\bigl(y\log p + (1-y)\log(1-p)\bigr).
\]

\textbf{Interprétation.}  
Si le modèle attribue une faible probabilité à la bonne classe, la pénalité est grande.

\subsection{Pourquoi cette loss ?}

Parce qu’elle est bien adaptée aux probabilités et qu’elle punit fortement les erreurs confiantes.

\section{Transfer learning}

\subsection{Idée générale}

On réutilise un réseau déjà entraîné sur une grande base d’images,
par exemple ImageNet.

On écrit conceptuellement :

\[
f(x)=g(h(x)),
\]

où :
\begin{itemize}
    \item $h(x)$ extrait des caractéristiques générales ;
    \item $g(x)$ est la tête de classification adaptée à notre tâche.
\end{itemize}

\subsection{Pourquoi cela marche ?}

Les premières couches apprennent des structures génériques :
\begin{itemize}
    \item bords ;
    \item textures ;
    \item motifs élémentaires.
\end{itemize}

Ces motifs sont utiles dans beaucoup de tâches visuelles, y compris en imagerie médicale.

\section{Évaluation du modèle}

\subsection{Matrice de confusion}

Pour une classification binaire, on distingue :
\begin{itemize}
    \item vrais positifs (TP),
    \item faux positifs (FP),
    \item vrais négatifs (TN),
    \item faux négatifs (FN).
\end{itemize}

\subsection{Accuracy}

\[
\mathrm{Accuracy}
=
\frac{TP+TN}{TP+TN+FP+FN}.
\]

\subsection{Précision}

\[
\mathrm{Precision}
=
\frac{TP}{TP+FP}.
\]

\subsection{Rappel}

\[
\mathrm{Recall}
=
\frac{TP}{TP+FN}.
\]

\subsection{Score F1}

\[
F_1 = 2\frac{\mathrm{Precision}\cdot \mathrm{Recall}}
{\mathrm{Precision}+\mathrm{Recall}}.
\]

\textbf{Pourquoi le F1 est utile ?}  
Parce qu’il pénalise les modèles déséquilibrés : il faut être bon à la fois en précision et en rappel.

\section{Courbe ROC et AUC}

La courbe ROC trace :
\begin{itemize}
    \item le taux de vrais positifs ;
    \item en fonction du taux de faux positifs.
\end{itemize}

On définit :

\[
\mathrm{TPR}=\frac{TP}{TP+FN},
\qquad
\mathrm{FPR}=\frac{FP}{FP+TN}.
\]

L’aire sous la courbe ROC est l’AUC.

\textbf{Interprétation.}  
Plus l’AUC est proche de 1, meilleur est le classifieur.

\section{Explicabilité : Grad-CAM}

Grad-CAM permet de visualiser les régions de l’image qui influencent la décision.

Si $A^k$ désigne la $k$-ième carte de caractéristiques de la dernière couche convolutive, on calcule

\[
\alpha_k^c
=
\frac{1}{Z}
\sum_{i,j}
\frac{\partial y^c}{\partial A_{i,j}^k}.
\]

Puis :

\[
L_{\mathrm{GradCAM}}^c
=
\mathrm{ReLU}
\left(
\sum_k \alpha_k^c A^k
\right).
\]

\textbf{Idée.}  
On obtient une carte de chaleur montrant \emph{où} le réseau regarde pour prendre sa décision.

\section{Lien direct avec le projet MedVision}

Dans MedVision :
\begin{itemize}
    \item les images sont chargées puis normalisées ;
    \item elles sont redimensionnées pour entrer dans un backbone CNN ;
    \item un modèle pré-entraîné est adapté par transfer learning ;
    \item l’entraînement est suivi expérimentalement avec MLflow ;
    \item les performances sont évaluées par plusieurs métriques ;
    \item l’inférence peut être expliquée via des cartes d’attention de type Grad-CAM.
\end{itemize}

\section{Conclusion}

Le projet MedVision illustre comment des outils issus :
\begin{itemize}
    \item de l’algèbre linéaire ;
    \item de l’analyse ;
    \item des probabilités ;
    \item et de l’optimisation ;
\end{itemize}
permettent de construire un pipeline moderne de computer vision pour l’imagerie médicale.

L’idée fondamentale est toujours la même :
\begin{center}
\emph{transformer des images en représentations mathématiques pertinentes,
puis apprendre automatiquement une règle de décision à partir des données.}
\end{center}

\end{document}